\section{Method}
\label{sec:method}

We adopt Differential Item Functioning (DIF) analysis from educational measurement to assess whether individual benchmark items behave consistently across groups of LLMs. The procedure requires only a binary response matrix and group labels, with no access to model internals or training data. Below we define the data, the statistical test, the cohort designs, a semi-synthetic validation framework, diagnostic checks, and statistical controls.

\subsection{Data}
\label{sec:data}

The primary dataset is the METABENCH response matrix~\citep{kipnis2024metabench}, which records binary outcomes (1~=~correct, 0~=~incorrect) for 5,227 models evaluated on 12,508 MMLU items~\citep{hendrycks2021mmlu}. Models span release dates from 2023 to 2024, drawn from the Open LLM Leaderboard. MMLU covers 57 subjects grouped into five broad domains (STEM, Humanities, Social Sciences, Other, and a mixed category). We supplement the response matrix with the web-overlap contamination labels of \citet{li2024contamination}, which classify 9,229 of 12,508 items (73.8\% coverage) as ``contaminated'' or ``clean'' based on verbatim matches in publicly accessible web corpora. These labels serve as an external reference for evaluating the relationship between DIF flags and existing contamination indicators.

\subsection{Mantel-Haenszel DIF and ETS Classification}
\label{sec:mhdif}

In educational testing, Differential Item Functioning (DIF) occurs when an item's statistical relationship to the latent trait $\theta$ differs between two groups after conditioning on ability level. An item with DIF is not merely harder for one group; it functions differently relative to what the group's overall ability would predict. The standard terminology refers to test-takers as ``examinees'' and questions as ``items.'' In our setting, examinees are LLMs and items are benchmark questions.

The Mantel-Haenszel (MH) procedure~\citep{holland1988mh} detects DIF by stratifying examinees into $K$ score bins based on total score (a proxy for $\theta$), where $K$ equals the number of distinct total scores observed in the combined sample, and computing a common odds ratio across strata. Within each stratum $k$, the responses of the reference group and the focal group form a $2 \times 2$ contingency table:

\begin{equation}
\label{eq:mh}
\alpha_{\mathrm{MH}} = \frac{\displaystyle\sum_{k=1}^{K} A_k D_k \,/\, T_k}{\displaystyle\sum_{k=1}^{K} B_k C_k \,/\, T_k}\,,
\end{equation}

\noindent where $A_k$ and $B_k$ are the number of correct and incorrect responses in the reference group for stratum $k$, $C_k$ and $D_k$ are the corresponding counts for the focal group, and $T_k$ is the total number of examinees in stratum $k$. Under the null hypothesis of no DIF, $\alpha_{\mathrm{MH}} = 1$. The log odds ratio is converted to the ETS delta scale:

\begin{equation}
\label{eq:delta}
\Delta_{\mathrm{MH}} = -2.35 \ln(\alpha_{\mathrm{MH}})\,.
\end{equation}

The Educational Testing Service (ETS) classifies DIF severity into three categories based on $|\Delta_{\mathrm{MH}}|$. Category~A (negligible) requires $|\Delta_{\mathrm{MH}}| < 1.0$. Category~C (large) requires both $|\Delta_{\mathrm{MH}}| \geq 1.5$ and statistical significance at $p < 0.05$ from the MH chi-squared test. Category~B (moderate) includes all remaining items: those with $1.0 \leq |\Delta_{\mathrm{MH}}| < 1.5$, or $|\Delta_{\mathrm{MH}}| \geq 1.5$ without statistical significance. The dual threshold in ETS~C (effect size \emph{and} significance) guards against flagging items that show large but unreliable effect sizes. Throughout this paper, we report the percentage of items classified as ETS~C, denoted C\%.

The choice of matching variable (the total score used for stratification) affects both validity and interpretation. We distinguish three matching strategies: (1)~\emph{standard matching} uses the within-benchmark total score on all items of the same benchmark; (2)~\emph{cross-subject matching} uses the total score on all subjects \emph{other than} the one containing the item under test, preventing the matching variable from absorbing injected contamination signals (\Cref{sec:injection}); and (3)~\emph{external matching} uses the total score on a different benchmark entirely (e.g., MMLU total score as the matching variable when testing DIF on GSM8K items), providing an independent ability proxy uncontaminated by the target benchmark. Cross-subject and external matching both address matching variable contamination but at different granularities: cross-subject operates within a multi-domain benchmark, while external operates across benchmarks.

\subsection{Cohort Designs}
\label{sec:cohorts}

DIF requires defining a reference group and a focal group. We construct five cohort comparisons, each designed to isolate a different source of variation.

\paragraph{Temporal.} The reference group consists of models released in 2023 ($N = 1{,}080$) and the focal group consists of models released in 2024 ($N = 807$). This is the primary comparison for detecting items whose functioning changed across model generations.

\paragraph{Ability-quintile.} Models within the same temporal window are divided by total score into quintiles. The focal group contains the lowest-scoring models (Q1--Q2) and the reference group contains the highest-scoring models (Q4--Q5). This comparison tests whether ability differences alone produce DIF, without temporal confounds.

\paragraph{Within-family.} Llama-2 models ($N = 161$) serve as the reference group and Llama-3 models ($N = 308$) as the focal group. This comparison isolates generational changes within a single model family.

\paragraph{Within-subject.} DIF is computed separately for each of the 57 MMLU subjects, with the matching variable (total score) restricted to items within the same subject. This design tests whether instability concentrates in specific knowledge domains.

\paragraph{Random split.} The 2023 cohort is divided into two equal random halves ($N = 540$ each). This serves as a sanity check: no systematic DIF should appear between two random subsets of the same population (expected C\% $\approx 0$).

\subsection{Semi-Synthetic Injection Framework}
\label{sec:injection}

To evaluate whether MH-DIF can detect differential contamination when it is present, we construct a semi-synthetic framework with controlled ground truth.

\paragraph{Procedure.} From items classified as ETS~A in the uninjected data (negligible DIF), we select 20\% uniformly at random as the ``contaminated'' set. For each contaminated item, we flip incorrect responses to correct in the focal group at exploitation rate $p \in \{0.0, 0.3, 0.5, 0.7, 1.0\}$. We then run MH-DIF on the modified response matrix and measure detection performance. We report $\Delta\mathrm{TPR} = \mathrm{TPR}(\text{injected}) - \mathrm{TPR}(\text{baseline})$ and $\Delta\mathrm{FPR} = \mathrm{FPR}(\text{injected}) - \mathrm{FPR}(\text{baseline})$, where TPR is the true positive rate on injected items and FPR is the false positive rate on uninjected items. Each injection level is repeated across 5 random seeds for uncertainty estimation.

\paragraph{Cross-subject matching.} The choice of matching variable is critical for valid semi-synthetic evaluation. When contamination is injected into items within subject $X$, the total score on subject $X$ absorbs part of the injection signal. Using this contaminated total score as the MH matching variable violates the conditional independence assumption and inflates FPR. Cross-subject matching resolves this problem by computing the matching variable from all subjects \emph{other than} the one being tested. Under this protocol, $\Delta\mathrm{FPR} = 0.000$ across all five MMLU domains, while detection power is preserved ($\Delta\mathrm{TPR} > 0.3$ at $p = 0.5$ for four of five subjects). We adopt cross-subject matching as the standard protocol for all semi-synthetic analyses.

\subsection{Diagnostics}
\label{sec:diagnostics}

Three diagnostic procedures complement the primary DIF analysis.

\paragraph{Local independence (Yen's Q3).} The MH procedure assumes that item responses are conditionally independent given $\theta$. Yen's Q3 statistic measures residual correlations between item pairs after conditioning on the latent trait. We compute Q3 for all within-subject and cross-subject item pairs, with $|Q3| > 0.20$ as the threshold for local dependence. Elevated within-subject dependence would indicate that items within the same subject share variance beyond what $\theta$ captures, which could inflate within-subject DIF rates.

\paragraph{Domain-specificity.} We test whether temporal DIF concentrates in domains where models improved most. At the ecological level, we compute Spearman $\rho$ between subject-level C\% and subject-level accuracy improvement from 2023 to 2024. At the item level, logistic regression predicts item-level DIF~C classification from domain-level accuracy change, with and without item property controls (difficulty, discrimination). This decomposition quantifies how much of the observed instability is attributable to aggregate domain-level gains versus item-specific behavioral change.

\paragraph{Instability decomposition.} When domain accuracy improvement is included alongside item-level controls, we compute the counterfactual C\% that would remain after removing domain-level effects. The gap between observed and counterfactual C\% measures the contribution of domain-level improvement to the total instability.

\subsection{Statistical Considerations}
\label{sec:statistical}

\paragraph{Multiple testing.} The ETS~C classification imposes a dual threshold (effect size $|\Delta_{\mathrm{MH}}| \geq 1.5$ and significance $p < 0.05$), which provides implicit control over false positives. As a robustness check, we apply Benjamini-Hochberg FDR correction at $q = 0.05$. The random-split sanity check provides an empirical estimate of the false positive rate under the null hypothesis.

\paragraph{Threshold sensitivity.} We report results across a range of DIF thresholds ($|\Delta_{\mathrm{MH}}| \in \{1.0, 1.5, 2.0, 2.5\}$) to verify that qualitative conclusions do not depend on the specific ETS~C cutoff. The standard threshold of $|\Delta_{\mathrm{MH}}| \geq 1.5$ corresponds to the ETS~C definition used throughout the paper.

\paragraph{Architecture heterogeneity.} The METABENCH MMLU population consists of 99.98\% decoder-only models (5,226 of 5,227). This near-complete architectural homogeneity eliminates architecture as a confound for DIF analysis. We verify this by excluding the single encoder-decoder model and confirming identical results.
